\section{Tightness}
\label{sec.tightness}

The quantity $\len{\ch}-\len{\ch\meet\ch'}$ is the number of blocks of $\ch$
that are not on $\ch'$. A client holding $\ch$ must prune at least this many
blocks to obtain a prefix of $\ch'$. \Cref{thm.prefix} bounds it by
$\frac{\advrtwo}{1-\advrtwo}(\len{\ch'}-\len{\ch})$ plus a constant. This
section gives an adversary against which the coefficient is attained up to a
factor that tends to $1$ as $f$ and $\epsilon$ tend to $0$: for every
$c<t/(n-2t)$ and every $b\ge0$, in a long enough execution and except with
small probability, $\len{\ch}-\len{\ch\meet\ch'}>c(\len{\ch'}-\len{\ch})+b$
for the chain $\ch$ sent to the client and every chain $\ch'$ adopted by an
honest validator. \Cref{cor.roundtight} shows the same for the pruning
depth of \Cref{cor.round}. We make no claim about the additive terms.

\begin{definition}[Private mining]\label{def.private}
    The strategy $\pipriv$ with parameters $r_0\ge1$ and $s\ge1$ is the
    following. Write $r=r_0+s$.
    \begin{enumerate}
        \item In every round $i$ of $[1,r_0]$, after the honest validators
            have sent their messages of round $i$, the adversary spends every
            query on extending a longest chain among the genesis chain and
            the chains sent by any validator at rounds up to $i$. Each block
            it computes extends the block it computed before in the same
            round. It sends every block at once.
        \item Let $\ch_0$ be a longest chain adopted by an honest validator at
            round $r_0+1$.
        \item In every round of $[r_0+1,r]$ the adversary spends every query
            on extending $\ch_0$. The first block it computes extends
            $\head{\ch_0}$, and every later block extends the block it
            computed before. It sends nothing to the validators.
        \item At round $r$ it sends to the client the chain $\ch$ formed by
            $\ch_0$ and all the blocks it computed in the rounds of
            $[r_0+1,r]$.
    \end{enumerate}
\end{definition}

A block computed at round $i$ may be extended by a query made at round $i$.
\Cref{lem.length} counts the adversarial blocks by $Z$ and already allows
this. Throughout this section, fix $r_0\ge1$ and $s\ge1$, write $r=r_0+s$,
$S_1=[1,r_0]$ and $S_2=[r_0+1,r]$, let $\ch_0$ and $\ch$ be as in
\Cref{def.private}, and write $\ell_0=\len{\ch_0}$.

\begin{lemma}\label{lem.attack}
    In a typical execution against $\pipriv$, the following hold.
    \begin{enumerate}
        \item $\ell_0=1+X(S_1)+Z(S_1)$, and every chain sent by round $r_0$
            has length at most $\ell_0$.
        \item $\ch\in\bt_r$ and $\len{\ch}=\ell_0+Z(S_2)$.
        \item Every chain $\ch'$ adopted by an honest validator at round $r$
            satisfies $\len{\ch}-\len{\ch\meet\ch'}\ge Z(S_2)$ and
            \[
                X(S_2)-Z(S_2)-1 \le \len{\ch'}-\len{\ch} \le X(S_2)-Z(S_2) .
            \]
    \end{enumerate}
\end{lemma}
\begin{proof}
    For the first item we show by induction on $i\in[1,r_0+1]$ that every
    honest validator adopts at round $i$ a chain of length
    $L_i=1+X([1,i-1])+Z([1,i-1])$, and that the longest chain among the
    genesis chain and the chains sent by round $i-1$ has length $L_i$. Both
    hold at round $1$. Assume both at a round $i\le r_0$. Honest validators send at round $i$ chains of length at most $L_i+X_i$,
    and one of length $L_i+1$ if $X_i=1$. So, before the adversary acts, the
    longest chain among the genesis chain and the chains sent at rounds up to
    $i$ has length $L_i+X_i$. The adversary
    extends it by $Z_i$ blocks and sends the result, of length $L_{i+1}$. At
    round $i+1$ every honest validator receives it, and the chain it has is
    at most $L_i+X_i$ long, so it adopts a chain of length $L_{i+1}$. The
    first item is the case $i=r_0+1$.

    In the rounds of $S_2$ the adversary makes queries only on the private
    branch, so $\len{\ch}=\ell_0+Z(S_2)$, and every block of $\ch$ was
    computed by round $r$, so $\ch\in\bt_r$.

    For the third item, the blocks the adversary computes in the rounds of
    $S_2$ are sent to the client only, so none of them lies on $\ch'$. Hence
    $\ch\meet\ch'\preceq\ch_0$ and
    $\len{\ch}-\len{\ch\meet\ch'}\ge\len{\ch}-\ell_0=Z(S_2)$. The blocks of
    $\ch'$ computed in rounds of $S_2$ are honest and, by \Cref{lem.honestone},
    at most $X(S_2)$ in number. The other blocks of $\ch'$ form a prefix of
    $\ch'$, since no predictions occur. This prefix is the genesis chain or a prefix of a chain sent at the round
    its tip was computed, so it has length at most
    $\ell_0$ by the first item. Hence $\len{\ch'}\le\ell_0+X(S_2)$. For the
    lower bound, if $s=1$ then $\len{\ch'}=\ell_0$ by the first item and
    $X(S_2)\le1$. If $s\ge2$, some honest validator has a chain of length at
    least $\ell_0+X_{r_0+1}$ at round $r_0+2$, and \Cref{lem.growth} with
    $v=r_0+2$ and $w=r$ gives
    $\len{\ch'}\ge\ell_0+X([r_0+1,r-1])\ge\ell_0+X(S_2)-1$.
\end{proof}

\begin{theorem}[Private mining]\label{thm.attack}
    Let $\epsilon'\in(0,1)$ and consider an execution against $\pipriv$. Except
    with probability at most $\typ+4e^{-\epsilon'^2\advr fs/3}$, every chain
    $\ch'$ adopted by an honest validator at round $r$ satisfies
    \begin{gather}
        \len{\ch}-\len{\ch\meet\ch'} > (1-\epsilon')\advr fs,
            \label{eq.attfork}\\
        \len{\ch'}-\len{\ch} < \left[(1+\epsilon')-(1-\epsilon')\advr\right]fs,
            \label{eq.attupper}\\
        \len{\ch'}-\len{\ch} > \left[(1-\epsilon')-(1+\epsilon')\frac{\advr}{1-f}\right]fs-1.
            \label{eq.attlower}
    \end{gather}
    Except with a further probability at most $4e^{-\epsilon'^2\advr fr_0/3}$,
    moreover
    \begin{equation}\label{eq.attheight}
        \begin{gathered}
            (1-\epsilon')\left[(1+\advr)fr_0+\advr fs\right] < \len{\ch}-1 \\
            < (1+\epsilon')\left[\Big(1+\frac{\advr}{1-f}\Big)fr_0+\frac{\advr}{1-f}fs\right].
        \end{gathered}
    \end{equation}
\end{theorem}
\begin{proof}
    $X(S_1)$ and $X(S_2)$ are sums of $r_0$ and $s$ independent Bernoulli
    trials with success probability $f$. $Z(S_1)$ and $Z(S_2)$ are sums of
    $tqr_0$ and $tqs$ independent Bernoulli trials with success probability
    $p$, and \Cref{eq.rates} gives $\advr fr_0<pqtr_0<\frac{\advr}{1-f}fr_0$
    and $\advr fs<pqts<\frac{\advr}{1-f}fs$. By the Chernoff bounds and
    $\advr<1$, each of the events
    \[
        \begin{gathered}
            X(S_2)\le(1-\epsilon')fs,\qquad X(S_2)\ge(1+\epsilon')fs,\\
            Z(S_2)\le(1-\epsilon')\advr fs,\qquad
            Z(S_2)\ge(1+\epsilon')\tfrac{\advr}{1-f}fs
        \end{gathered}
    \]
    has probability at most $e^{-\epsilon'^2\advr fs/3}$, and each of the four
    corresponding events for $X(S_1)$ and $Z(S_1)$, with $r_0$ in place of
    $s$, has probability at most $e^{-\epsilon'^2\advr fr_0/3}$. By a union
    bound, except with probability $\typ+4e^{-\epsilon'^2\advr fs/3}$ the
    execution is typical and none of the first four events occurs, and then
    the third item of \Cref{lem.attack} gives
    \Cref{eq.attfork,eq.attupper,eq.attlower}. If moreover none of the other
    four events occurs, the first two items of \Cref{lem.attack} give
    \Cref{eq.attheight}, since $\len{\ch}-1=X(S_1)+Z(S_1)+Z(S_2)$.
\end{proof}

\begin{corollary}[Tightness]\label{cor.tight}
    Let $0\le c<t/(n-2t)$ and $b\ge0$. There are $\epsilon'\in(0,1)$ and
    $s_0$, depending only on $c$, $b$, $\advr$ and $f$, such that for every
    $r_0\ge1$ and $s\ge s_0$, in an execution against $\pipriv$, except with probability at most $\typ+4e^{-\epsilon'^2\advr fs/3}$, every chain
    $\ch'$ adopted by an honest validator at round $r$ satisfies
    $\len{\ch'}>\len{\ch}$ and
    \[
        \len{\ch}-\len{\ch\meet\ch'} > c\left(\len{\ch'}-\len{\ch}\right)+b .
    \]
\end{corollary}
\begin{proof}
    Since $c<\advr/(1-\advr)$, we have $\advr(1+c)>c$. From \Cref{lem.rho},
    $\advr/(1-f)<1$. Choose $\epsilon'\in(0,1)$ small enough that
    \[
        \delta_1=(1-\epsilon')\advr(1+c)-c(1+\epsilon')>0
        \quad\text{and}\quad
        \delta_2=(1-\epsilon')-(1+\epsilon')\frac{\advr}{1-f}>0 .
    \]
    Consider an execution in which
    \Cref{eq.attfork,eq.attupper,eq.attlower} hold. From
    \Cref{eq.attfork,eq.attupper},
    \[
        \len{\ch}-\len{\ch\meet\ch'}-c\left(\len{\ch'}-\len{\ch}\right)
        > (1-\epsilon')\advr fs-c\left[(1+\epsilon')-(1-\epsilon')\advr\right]fs
        = \delta_1fs .
    \]
    From \Cref{eq.attlower}, $\len{\ch'}-\len{\ch}>\delta_2fs-1$. Take
    $s_0\ge\max\{b/(\delta_1f),\ 1/(\delta_2f)\}$. For $s\ge s_0$ the two
    right-hand sides are at least $b$ and $0$.
\end{proof}

\begin{corollary}[Tightness in $r$ and $\len{\ch}$]\label{cor.roundtight}
    Let $\epsilon'\in(0,1)$, $K\ge1$ and $r_0\le Ks$, and consider an
    execution against $\pipriv$ in which \Cref{eq.attfork,eq.attheight} hold.
    Let $\prn(r)$ be the maximum in \Cref{cor.round} for $\ch$ and $r$. Then
    $\len{\ch}-\len{\ch\meet\ch'}>(1-\epsilon')\advr fs$ for every chain
    $\ch'$ adopted by an honest validator at round $r$, and
    \begin{gather*}
        \prn(r)\le\max\left\{k_0,\ \frac{\advrone B+\maxr(2f+1)}{1-f_2/f}\right\},\\
        B=\left[1+(1+K)\slack+\epsilon'(\advr+(1+\advr)K)\right]fs+\maxr f\lambda .
    \end{gather*}
\end{corollary}
\begin{proof}
    The first claim is \Cref{eq.attfork}. Since $\maxr=1+\advr+\slack$,
    \Cref{eq.attheight} gives
    \[
        \maxr f(r+\lambda)-\len{\ch}+1
        < (\epsilon'(1+\advr)+\slack)fr_0+(1+\epsilon'\advr+\slack)fs+\maxr f\lambda
        \le B ,
    \]
    where the last step uses $r_0\le Ks$. The second term of the maximum in
    \Cref{cor.round} is
    $\frac{\advrone(\maxr f(r+\lambda)-\len{\ch}+1)+\maxr(2f+1)}{1-f_2/f}$.
\end{proof}

For large $s$ the pruning depth of \Cref{cor.round} therefore exceeds
$\len{\ch}-\len{\ch\meet\ch'}$ by a factor of at most
$\frac{\advrone[1+(1+K)\slack+\epsilon'(\advr+(1+\advr)K)]}{(1-\epsilon')\advr(1-f_2/f)}$
plus a term that tends to $0$ with $s$. This factor tends to $1$ as
$\epsilon'$, $f$ and $\epsilon$ tend to $0$. By \Cref{eq.attheight} the
height of $\ch$ is about $(1+\advr)fr-fs$, so every height between
$\advr fr$ and $(1+\advr)fr$ is obtained by a choice of $s$.
